%%% 前言.tex — 习题册前言 %%%

\chapter*{前言}
\addcontentsline{toc}{chapter}{前言}

如果你想要更改 PDF，请不要删改这里的前言和声明，谢谢。如果想了解更多，请访问我的主页\url{https://rikka.moe}

\medskip

本书的起点源自我在高中三年学习数学期间用 Obsidian 记下的笔记。写这本书的初衷很简单：我觉得市面上没有一本教辅可以真正满足我。

目前大部分教辅的编排方式是"总结+例题"的模式，实质上是让学生做 Pattern Matching——看到类似的题就套上去。这种模式对应试确实有效，但总觉得缺了什么。直到我读了 Springer 出版社的《初等数学之旅》，才意识到好的数学书应该展示清晰的路径图：解题需要追踪一道题背后的逻辑流，从哪个定义或定理出发，经历哪些逻辑分支，最后收敛到答案。

这也是为什么我在《数学之旅》系列中坚持"所有结论都要有证明"：每一个结论都必须伴随路径，没有路径的结论不应存进你的脑子里。而本书作为该系列的配套习题集，则从另一个角度出发——为每一条路径提供练习的机会。知识框架梳理思路，习题集锤炼手感，两者相辅相成。

本书按章节编排，每章以知识框架开篇，梳理核心概念与定理，随后分为选择题和解答题两组。习题册采用双栏排版：左栏为题目，右栏为答题区，方便学习者直接在书上作答。答案册与习题册分册发行，每道题均附详细解析。

本书目前为书稿状态，内容可能尚未完整，文中可能存在笔误或疏漏。若发现错误或有建议，欢迎联系作者。

\medskip

本版本为 2026 年 6 月 27 日的模块化重构版本，采用了全新的 xsim 题库架构，实现内容与格式的彻底分离。具体更新内容请参见 GitHub Release 页面，或访问本书落地页：\url{https://rikka.moe/books}。

\medskip

感谢你阅读本书。

\bigskip

本书属于 Project Akari 系列。Akari 取自日文「あかり」，意为光，寓意是让好的学习材料更容易被找到、更容易被用起来。

Project Akari 的另一部书稿《初等数学之旅》将中学数学内容系统展开，每章从概念到定理、从证明到应用，展示完整的逻辑路径。两本书可以独立阅读，也可以互相参照——一册讲"为什么"，一册练"怎么做"。配套的 Akari Web 平台（\url{https://web.chronorise.com}）提供笔记、记忆卡片和在线练习，与书稿配合使用。

本系列的基础版本免费开放，完整版本以极低价格提供。知识不应该是奢侈品，这是我做这件事的底线。

从写作、排版到软件开发，Project Akari 的每一步都是我一个人完成的。服务器和运营成本靠书的销售收入以及我个人的工作收入来覆盖。这个项目既是学习方法的实践，也是工程能力的验证。

\cleardoublepage
